\resumeSubheadingReza
{Automotive Software Engineer}{05/2022 - 05/2023}
{\rezaLink{https://sw-motion.tech/}{SoftwareMotion}}{Remote, China}
% {01/2023 - Now}{Remote, China}
{
  % At \rezaLink{https://sw-motion.tech/}{SoftwareMotion}, developed a \textbf{CI/CD} pipeline in \textbf{Gitlab} using \textbf{GTest} for the Autosar-based codebase, enhancing testing efficiency by 30\%, and created a \textbf{PyQT} visualization tool for AEB and ACC systems, significantly reducing debugging time by processing MF4 and CAN signals.
  % SoftwareMotion is a company that develops Advanced Driver Assistance Systems(ADAS). I worked there as an Algorithm Engineer in the Decision Planning and Control(DPC) team. My major projects was to implement a CI procedure in gitlab and making a visualization app for debugging our functionalities. These tools reduced debugging time significantly.
  
  \vspace{-12pt}
    \begin{itemize}
      \item {Evaluated and integrated advanced path planning solutions, such as lattice-based and \rezaLink{https://github.com/ApolloAuto/apollo}{Apollo}'s EM planner. Also, developed obstacle avoidance algorithms with \textbf{DP} and \textbf{A*} techniques, validated through a custom Python-based simulator.}
      \item {Made a \textbf{CI/CD} pipeline in \textbf{Gitlab} with \textbf{GTest} for the Autosar-based codebase(\textbf{C/C++}), reducing on-field testing through automated unit and integration tests by 30\%}
      \item {Built a \textbf{PyQT} visualization software for AEB and ACC systems, integrating BEV plots to drastically reduce debugging time. Processed MF4 and CAN signals (BLF, DBC) to expedite vehicle issue diagnosis.}
  \end{itemize}
}